\chapter{Conclusions and Future Work}

This final chapter summarizes the main contributions and methodological advances achieved throughout the thesis and outlines recommended directions for future work and technology transfer.

\section*{Conclusions}

This thesis systematically advances the state of the art in both linear and nonlinear projection-based reduced order modeling by developing, implementing, and rigorously assessing new methodologies across a range of test cases, from canonical model problems to realistic industrial configurations. Key contributions include the formulation of robust hyper-reduction strategies for Petrov–Galerkin PROMs, the design of scalable HPC-enabled workflows for intrusive ROM training, and the development of novel nonlinear manifold approaches, such as PROM-ANN, PROM-GPR, and PROM-RBF, that mitigate the Kolmogorov $n$-width barrier.

A significant practical outcome is the extension of two open-source frameworks, Kratos Multiphysics and AERO-F, to support these advanced ROM techniques, enabling systematic testing from simplified problems like the Burgers equation to complex scenarios including nonlinear structural mechanics, incompressible Navier–Stokes flows, compressible aerodynamic benchmarks, and multiphysics cases such as the thermal dynamics of an electric motor in collaboration with Siemens and the Ahmed body wake flow problem in collaboration with the Farhat Research Group at Stanford University. This validation across diverse physical phenomena illustrates both the generality and robustness of the proposed methodologies.

While this thesis primarily focused on the development and rigorous evaluation of projection-based ROMs, it is important to note that no full digital twin was implemented. The PROMs developed here remain on the testing side and have not yet been connected to live physical assets, a prerequisite for realizing the closed-loop functionality of a true digital twin. Accordingly, the work presented can be seen as laying the foundation for future \textit{Component Twins} and \textit{Asset Twins}: digital representations of individual parts or complete physical assets that can be integrated into broader digital twin ecosystems.

Lastly, the impact of this work extends beyond academic contributions by planting a seed for technology transfer through the \textit{SimTwins} spin-off initiative. This effort aims to leverage the open-source capabilities of projection-based ROMs to deliver scalable digital shadows and digital twins that meet the emerging demands of Industry 4.0 and 5.0, thus positioning projection-based ROM techniques as central enabling technologies in the current and future deep-tech landscape.

\section*{Future Work}

Building on the findings and methodologies established in this thesis, several avenues are recommended for future research and development:
\begin{itemize}
    \item \textbf{Piecewise nonlinear closure modeling:} Extend nonlinear latent-space closure strategies (PROM-ANN, PROM-RBF, PROM-GPR) to piecewise or local subspace frameworks, leveraging Local POD or clustering strategies. This will further enhance efficiency and accuracy in complex, strongly nonlinear, or convection-dominated industrial problems.
    
    \item \textbf{Physics-embedded machine learning:} Further refine and expand methods that embed governing physics into machine learning-based ROMs, including residual-consistent training (such as the discrete physics-informed strategy demonstrated here) and projection-based hybrid formulations. Prioritize such approaches in settings where governing equations are available, thereby strengthening the reliability and robustness of forward simulations in industrial contexts.
    
    \item \textbf{Broader industrial application:} Expand these ROM techniques into increasingly challenging industrial configurations, particularly hypersonic flows and multiphysics models with extensive parametric spaces. Such applications will fully leverage and demonstrate the scalability and computational efficiency achievable by projection-based ROMs.
    
    \item \textbf{Closing the digital twin loop:} Focus efforts on integrating developed PROMs with real-time sensor data and physical assets, transitioning from robustly tested surrogates into operational \textit{Component Twins} and \textit{Asset Twins}. This integration is essential for continuous monitoring, adaptation, and real-time decision support capabilities in practical engineering systems.
    
    \item \textbf{Support for open-source frameworks and technology transfer:} Continue enhancing, maintaining, and contributing to open-source frameworks such as Kratos Multiphysics and AERO-F, ensuring advanced ROM techniques remain widely accessible. Simultaneously, further development through the \textit{SimTwins} initiative will be crucial for translating research advances into practical and industry-ready ROM-based solutions.
\end{itemize}

Together, these future directions aim to consolidate the impact of this thesis and further position projection-based reduced order modeling as a robust, scalable, and physically consistent tool for real-time digital twins and next-generation engineering applications.
